\chapter{联邦学习}
\label{chap:federated}
\begin{quote}\small
\textit{本章摘要。}联邦学习在原始数据由多个客户端分别持有时组织协作训练。它把统计目标、客户端抽样、本地计算、通信协议和隐私威胁放到同一个学习流程中。本章先区分横向联邦中的客户端、样本权重和参与规则，再由集中式梯度下降推导联邦平均（Federated Averaging，FedAvg）。随后用客户端漂移解释 FedProx、SCAFFOLD 和 FedOpt 的机制差异，介绍个性化联邦学习、安全聚合与差分隐私。最后给出评价协议和失败情形。\end{quote}

上一章的主动学习假设标注样本能够进入同一训练集。现实中的工厂、医院或移动设备可能受权限、带宽和法规限制，不能把原始记录上传到共同服务器。联邦学习保留这一数据边界，让客户端在本地计算模型更新，服务器只协调模型或更新。它改变了数据的流动方式，但没有改变预测任务本身。

例如，三家工厂都希望预测轴承故障。每家工厂可以用自己的振动记录更新同一个模型，却不能交换原始波形。工厂的数据量、故障比例、设备性能和在线时间可能不同。服务器还要处理掉线和恶意更新。由此产生四个需要同时回答的问题。全局模型究竟代表哪一类样本。多步本地训练为什么会偏离全局目标。服务器能看到哪些信息。平均准确率是否掩盖了某个工厂的失效。

\section{协作边界与基本对象}

联邦学习适合多个参与方执行相同预测任务，且各方能够使用相同的模型结构和参数坐标。横向联邦（horizontal federated learning）中，不同客户端拥有相似特征和标签空间，但样本不同。例如不同工厂都记录温度、振动和故障类别。纵向联邦（vertical federated learning）中，不同机构拥有同一批对象的不同特征。纵向场景需要样本对齐、加密计算或安全求交，不能直接平均参数。本章只讨论横向联邦。

\begin{figure}[htbp]
\centering
\begin{tikzpicture}[node distance=0.75cm and 0.9cm]
\node[idi state,minimum width=2.4cm] (server) {协调服务器\\全局参数 $\theta^t$};
\node[idi input,below left=1.0cm and 1.4cm of server,minimum width=2.1cm] (a) {客户端 $1$\\本地数据 $D_1$};
\node[idi input,below=1.0cm of server,minimum width=2.1cm] (b) {客户端 $2$\\本地数据 $D_2$};
\node[idi input,below right=1.0cm and 1.4cm of server,minimum width=2.1cm] (c) {客户端 $K$\\本地数据 $D_K$};
\draw[idi arrow] (server)--(a) node[midway,left,font=\scriptsize]{下发模型};
\draw[idi arrow] (server)--(b);
\draw[idi arrow] (server)--(c);
\draw[idi arrow] (a)--(server) node[midway,left,font=\scriptsize]{上传更新};
\draw[idi arrow] (b)--(server);
\draw[idi arrow] (c)--(server);
\node[below=0.25cm of b,font=\scriptsize] {原始数据留在客户端，服务器聚合更新};
\end{tikzpicture}
\caption{横向联邦学习的一轮同步通信。服务器下发同一版本的全局参数，客户端本地计算后上传更新，服务器再按规定规则聚合。图中箭头表示消息方向，不表示原始样本流向。}
\label{fig:federated-basic-loop}
\end{figure}

图\ref{fig:federated-basic-loop}中的客户端可以是手机、工厂或医院。cross-device 联邦通常包含大量设备，每台设备数据少、在线时间短，系统必须容忍高掉线率。cross-silo 联邦通常只有少量组织，每个组织数据较多，参与方身份和网络连接相对稳定。两种设定都可以使用相同的优化公式，但客户端抽样、可信关系和通信预算不同。

设客户端 $k$ 持有数据集 $D_k=\{(x_{ki},y_{ki})\}_{i=1}^{n_k}$，其中 $n_k\geq1$ 是样本数。模型参数为 $\theta\in\mathbb R^d$。本地平均损失定义为
\begin{equation}
 F_k(\theta)=\frac{1}{n_k}\sum_{i=1}^{n_k}\ell(\theta;x_{ki},y_{ki}).
\end{equation}
这里 $\ell$ 是单样本损失，$d$ 是参数个数。服务器若希望每条样本在总体中等权，应取
\begin{equation}
 F(\theta)=\sum_{k=1}^{K}p_kF_k(\theta),\qquad
 p_k=\frac{n_k}{\sum_{j=1}^{K}n_j}.
 \label{eq:fed-global-objective-rewrite}
\end{equation}
若业务目标要求每家工厂等权，则取 $p_k=1/K$。$p_k$ 是目标权重，不是客户端本轮是否上线的概率。改变 $p_k$ 就改变了要优化的总体。样本量大的客户端在第一种目标中影响更大，小样本但高风险的客户端可能需要单独的代价权重。

客户端权重和样本权重还要与标签损失区分。若故障漏报比误报代价高，应在 $\ell$ 中表达类别代价。只增加某个工厂的 $p_k$，不能准确表示该工厂所有类别都更重要。设备数量也不是客户端数量。一个客户端可以管理许多设备，一个设备也可能被拆成多个逻辑客户端。划分方式应与权限、通信和测试集划分保持一致。

\section{FedAvg的目标与计算}

联邦学习最早受到分布式移动设备训练的推动。McMahan 等人提出的 FedAvg 将多步本地随机梯度下降与服务器加权平均结合起来\cite{mcmahan2017fedavg}。Kairouz 等人的综述系统整理了统计异质性、系统约束和隐私问题\cite{kairouz2019advances}。Li 等人随后从优化、系统和隐私三个方面总结了联邦学习的主要挑战\cite{li2020federatedsurvey}。近年来的综述进一步转向大语言模型和基础模型，Yao 等人讨论了联邦大语言模型的训练、参数高效适配和开放问题\cite{yao2024federatedllm}，Jiang 等人集中讨论可行性、鲁棒性和安全边界\cite{jiang2025federatedllm}。该设计的直接目的，是用更多本地计算换取更少的同步通信。它的代价是客户端在不同数据分布上移动，聚合更新可能偏离全局梯度。

先考虑所有客户端参与、确定性梯度和一步更新。客户端从同一个参数 $\theta^t$ 出发，计算
\begin{equation}
 \theta_k^{t+1}=\theta^t-\eta\nabla F_k(\theta^t),
 \qquad \eta>0.
\end{equation}
服务器按目标权重平均，得到
\begin{align}
 \sum_kp_k\theta_k^{t+1}
 &=\sum_kp_k\left(\theta^t-\eta\nabla F_k(\theta^t)\right)\\
 &=\theta^t-\eta\sum_kp_k\nabla F_k(\theta^t)\\
 &=\theta^t-\eta\nabla F(\theta^t).
 \label{eq:federated-centralized-equivalence}
\end{align}
式\eqref{eq:federated-centralized-equivalence}说明这一特殊情形与集中式全梯度下降完全相同。它需要所有客户端从相同版本开始，权重和为一，并且本地梯度确实对应式\eqref{eq:fed-global-objective-rewrite}。随机小批量梯度只在条件无偏时具有相同的期望，实际更新仍有方差。

FedAvg 将本地更新扩展为 $E$ 步。令 $\theta_{k,0}^t=\theta^t$，并令 $g_{k,e}^t$ 是客户端 $k$ 在本地点的随机小批量梯度，则
\begin{equation}
 \theta_{k,e+1}^t=\theta_{k,e}^t-\eta g_{k,e}^t,
 \qquad e=0,\ldots,E-1.
 \label{eq:fedavg-local-update-rewrite}
\end{equation}
本地计算结束后，服务器执行
\begin{equation}
 \theta^{t+1}=\sum_{k=1}^{K}p_k\theta_{k,E}^t.
 \label{eq:fedavg-aggregation-rewrite}
\end{equation}
也可以上传差分 $u_k^t=\theta_{k,E}^t-\theta^t$，然后使用 $\theta^{t+1}=\theta^t+\sum_kp_ku_k^t$。完整参数和差分在权重和为一时给出相同结果，但差分更容易配合裁剪、量化和安全聚合。

\paragraph{数值例子。}
两个客户端分别有 $1$ 条和 $3$ 条样本。服务器采用样本权重 $p_1=1/4$、$p_2=3/4$。从标量参数 $\theta^t=0$ 出发，本地梯度分别为 $2$ 和 $-2$，步长为 $0.1$。一步更新后，客户端参数为 $-0.2$ 和 $0.2$。服务器得到
\[
 \theta^{t+1}=\tfrac14(-0.2)+\tfrac34(0.2)=0.1.
\]
集中式梯度为 $\frac14(2)+\frac34(-2)=-1$，所以集中式一步更新也是 $0.1$。若改成客户端等权，聚合结果变为 $0$。这不是计算误差，而是优化目标发生了变化。

本地步数 $E$ 表示梯度更新次数，不等于数据集遍历次数。一个本地 epoch 可能包含多个小批量。增大 $E$ 可以减少同步频率，却可能增加客户端漂移。通信成本应记录下行和上行参数字节数，计算成本应记录本地步数和端侧时间，不能只比较通信轮数。

\section{部分参与与非独立同分布}

cross-device 系统通常每轮只抽取部分客户端。令 $S_t$ 为第 $t$ 轮参与集合。常用聚合为
\begin{equation}
 \theta^{t+1}=\sum_{k\in S_t}\widetilde p_k^t\theta_{k,E}^t,
 \qquad
 \widetilde p_k^t=\frac{p_k}{\sum_{j\in S_t}p_j}.
 \label{eq:federated-inclusion-gradient}
\end{equation}
该规则保证本轮权重和为一，但一般不再是原总体目标的无偏梯度。若上线概率为 $\pi_k>0$ 且可估计，可以使用逆概率校正 $p_k/\pi_k$。小概率客户端的更新会放大估计方差。抽样与网络状态相关时，简单增加轮数也不能恢复缺失群体的代表性。

非独立同分布（non-Independent and Identically Distributed，non-IID）表示客户端的联合分布 $P_k(X,Y)$ 不同。近期数据中心综述进一步把异质性、数据质量和收敛行为联系起来，说明客户端差异需要结合数据生成和训练过程共同分析\cite{nguyen2026datacentricfl}。标签比例不同属于 label skew。输入工况不同属于 feature shift。给定输入后标签规则不同属于 concept shift。客户端数据量不同属于 quantity skew。系统算力和网络速度不同属于系统异构，不等同于统计 non-IID，但两者会共同影响参与和收敛。

令全局梯度为 $\nabla F(\theta)=\sum_kp_k\nabla F_k(\theta)$。客户端间梯度差异可以用
\begin{equation}
 \zeta^2(\theta)=\sum_kp_k\|\nabla F_k(\theta)-\nabla F(\theta)\|_2^2
\end{equation}
描述。即使全局梯度为零，各个本地梯度也可能不为零。客户端各自继续训练时，参数会向不同的本地目标移动。这种由本地数据差异引起的偏离称为客户端漂移（client drift）。

在全参与和确定性梯度下，把式\eqref{eq:fedavg-local-update-rewrite}求和可得
\begin{align}
 \theta^{t+1}-\theta^t
 &=-\eta\sum_kp_k\sum_{e=0}^{E-1}\nabla F_k(\theta_{k,e}^t)\\
 &=-\eta E\nabla F(\theta^t)-\eta R_t,\\
 R_t&=\sum_kp_k\sum_{e=0}^{E-1}
 [\nabla F_k(\theta_{k,e}^t)-\nabla F_k(\theta^t)].
 \label{eq:fed-drift-decomposition-rewrite}
\end{align}
$R_t$ 是多步本地轨迹相对于固定起点梯度的偏离。若每个 $F_k$ 的梯度是 $L$-Lipschitz，且本地轨迹上的梯度范数不超过 $G$，则
\[
 \|\theta_{k,e}^t-\theta^t\|_2\leq\eta eG,
 \qquad
 \|R_t\|_2\leq L\eta G\frac{E(E-1)}2.
\]
这个上界只说明本地步数、步长和曲率会放大轨迹偏离。它不是一般非凸网络的完整收敛定理。

一个标量二次例子可以显示充分本地训练的风险。令
\[
 F_k(\theta)=\tfrac12a_k(\theta-b_k)^2,
 \qquad a_k>0.
\]
客户端训练 $E$ 步后的参数为
\[
 \theta_{k,E}=b_k+(1-\eta a_k)^E(\theta^t-b_k).
\]
若 $|1-\eta a_k|<1$ 且 $E$ 很大，聚合趋近于 $\sum_kp_kb_k$。而全局目标的最优点是 $\sum_kp_ka_kb_k/\sum_kp_ka_k$。例如等权取 $(a_1,b_1)=(1,0)$ 和 $(a_2,b_2)=(3,2)$，本地充分训练后的平均值为 $1$，全局最优点为 $1.5$。本地模型训练得更充分，未必使共享模型更接近全局目标。

\section{异质性下的优化方法}

FedAvg 是通信受限时的基线。数据分布接近、客户端数量大且每轮只能通信少量次数时，它通常容易实现。客户端漂移明显时，需要改变本地目标或更新方向。下面的方法按它们保存和使用的信息来区分。

\subsection{FedProx的近端约束}

FedProx 针对本地数据和设备差异，在每轮本地目标中加入服务器参数的近端惩罚\cite{li2020fedprox}：
\begin{equation}
 \min_v F_k(v)+\frac{\mu}{2}\|v-\theta^t\|_2^2,
 \qquad \mu\geq0.
 \label{eq:fedprox-objective-rewrite}
\end{equation}
$\mu$ 控制本地适应和全局一致之间的取舍。对应的梯度更新为
\begin{equation}
 v_{e+1}=v_e-\eta[\nabla F_k(v_e)+\mu(v_e-\theta^t)],
 \qquad v_0=\theta^t.
\end{equation}
第一步的近端梯度为零，因此与 FedAvg 相同。之后模型越远离 $\theta^t$，惩罚越大。$\mu$ 过小不能抑制漂移，过大则削弱本地信息。对非凸神经网络，近端项也不保证本地目标变成凸函数。

\subsection{SCAFFOLD的控制变量}

SCAFFOLD 用服务器控制变量 $c$ 和客户端控制变量 $c_k$ 估计系统性梯度偏移\cite{karimireddy2020scaffold}。本地更新改为
\begin{equation}
 v_{e+1}=v_e-\eta[g_k(v_e)-c_k+c].
 \label{eq:scaffold-update-rewrite}
\end{equation}
减去 $c_k$ 是去除客户端特有方向，加上 $c$ 是补回全局参考方向。若在同一参数点有 $c_k=\nabla F_k(\theta^t)$ 和 $c=\sum_kp_kc_k$，第一步更新的条件期望方向为全局梯度。后续参数点改变后，校正仍是近似。

控制变量需要随客户端保存，并可能在掉线后变得陈旧。它增加了内存和通信量，适合客户端能够长期保留状态的场景。系统若经常更换临时设备，应说明控制变量如何初始化和恢复。SCAFFOLD 改善的是优化漂移，不提供隐私保护。

\subsection{FedOpt的服务器优化}

FedOpt 将客户端聚合得到的平均差分看作服务器的梯度替代量，再使用服务器端优化器更新全局模型\cite{reddi2021adaptive}。令
\begin{equation}
 \Delta_t=\sum_{k\in S_t}\widetilde p_k^t(\theta_{k,E}^t-\theta^t).
\end{equation}
FedAvg 直接令 $\theta^{t+1}=\theta^t+\Delta_t$。FedOpt 则可以维护一阶动量 $m_t$ 和二阶统计量 $v_t$，例如 Adam 型服务器更新
\begin{align}
 m_t&=\beta_1m_{t-1}+(1-\beta_1)\Delta_t,\\
 v_t&=\beta_2v_{t-1}+(1-\beta_2)\Delta_t\odot\Delta_t,\\
 \theta^{t+1}&=\theta^t+\eta_s\frac{m_t}{\sqrt{v_t}+\varepsilon}.
 \end{align}
其中 $\odot$ 是逐坐标乘法，$\beta_1,\beta_2\in[0,1)$，$\eta_s>0$ 是服务器步长，$\varepsilon>0$ 防止除零。客户端差分不是严格梯度时，动量和自适应缩放可能改善或恶化训练。必须固定客户端步长、服务器步长和总通信预算后比较，不能把 FedOpt 的收益归因于聚合本身。

\begin{table}[htbp]
\centering
\caption{联邦优化方法的机制与代价}
\label{tab:federated-method-constraints}
\scriptsize
\begin{tabular}{m{1.7cm}m{2.5cm}m{2.7cm}m{2.5cm}}
\toprule
方法 & 主要解决的问题 & 保存或修改的信息 & 主要边界 \\
\midrule
FedAvg & 通信次数较多 & 客户端本地多步，服务器加权平均 & 多步和 non-IID 会产生漂移 \\
FedProx & 本地模型偏离过快 & 本地目标加入 $\|v-\theta^t\|^2$ & 近端系数需要验证，不能消除分布差异 \\
SCAFFOLD & 系统性客户端漂移 & 维护 $c$ 与 $c_k$ 控制变量 & 需要额外状态，掉线后可能陈旧 \\
FedOpt & 聚合更新的尺度和方向不稳定 & 服务器端动量或自适应二阶统计 & 超参数和差分噪声会改变行为 \\
\bottomrule
\end{tabular}
\end{table}

表\ref{tab:federated-method-constraints}按机制而不是按名称比较方法。若主要困难是通信，应先建立 FedAvg 基线。若本地数据和算力差异造成漂移，可比较 FedProx 或 SCAFFOLD。若服务器聚合差分的尺度变化明显，再考虑 FedOpt。每次比较都要说明改变了哪一个目标或更新步骤。

\section{个性化模型与基础模型时代}

共享模型假设客户端之间存在足够共同的预测规律。若不同工厂的传感器、工况或故障代价不同，一个模型可能无法同时达到各方要求。个性化联邦学习允许部分参数保留在客户端。常见做法是共享编码器 $f_\theta$，每个客户端保存预测头 $h_{\psi_k}$，并在本地优化 $h_{\psi_k}(f_\theta(x))$。只有结构和类别语义一致的共享参数才能逐坐标聚合。

另一种软共享目标为
\begin{equation}
 \min_{\theta,\{v_k\}}\sum_kp_k\left[F_k(v_k)+\frac{\lambda}{2}\|v_k-\theta\|_2^2\right],
 \qquad \lambda>0.
 \label{eq:personalized-federated-objective-rewrite}
\end{equation}
固定客户端模型 $v_k$ 时，最优共享参数满足 $\theta=\sum_kp_kv_k$。$\lambda$ 大时更接近共享模型，$\lambda$ 小时更接近独立模型。数据量少的客户端保留完整个性化模型可能过拟合，因此需要本地验证集和跨客户端测试。

预训练模型和基础模型改变了联邦训练的资源分配。近期联邦提示学习综述把提示参数视为共享模型与本地任务之间的中间层，强调参数高效微调、客户端差异和评测协议的共同影响\cite{yang2026federatedprompt}。大模型全参数上传的通信和隐私成本很高，实际系统常冻结大部分参数，只联邦训练低秩适配器、提示参数或任务头。低秩更新减少字节数，却可能泄露任务信息，也不自动解决不同客户端的标签语义差异。模型规模增大还会放大客户端算力差异。基础模型时代的联邦学习因此更关注参数高效微调、模型路由和个性化，而不是简单地把完整模型平均。

\section{隐私、通信与信任边界}

原始数据留在客户端只说明数据没有直接上传。梯度或参数差分仍可能包含样本信息。攻击者可能执行梯度反演、成员推断或属性推断。恶意客户端还可以上传投毒或后门更新。隐私、鲁棒性和公平性是不同目标，应分别声明威胁模型和评价方法。

安全聚合（secure aggregation）使服务器在协议完成后只能得到客户端更新的总和，而不能读取单个更新。一个简化的掩码说明如下。客户端 $k$ 与 $j$ 共享掩码 $r_{kj}$，并约定 $r_{jk}=-r_{kj}$。客户端发送
\begin{equation}
 \widetilde u_k=u_k+\sum_{j\ne k}r_{kj}.
\end{equation}
所有消息相加时，成对掩码抵消，得到 $\sum_k\widetilde u_k=\sum_ku_k$。真实协议还需要密钥协商、认证、门限恢复和掉线处理。该代数例子不是可部署的密码协议。

安全聚合隐藏单个更新，却不能阻止恶意客户端上传异常更新。它也不能防止服务器通过很小的聚合组或多次差分查询推断个体贡献。服务器若只能看到总和，也不能直接执行依赖单个更新的坐标中位数。安全聚合和鲁棒聚合需要联合设计。

差分隐私（Differential Privacy，DP）限制相邻数据集的输出分布差异。随机机制 $M$ 满足 $(\epsilon,\delta)$-差分隐私，是指任意相邻数据集 $D,D'$ 和事件 $A$ 都有
\begin{equation}
 \Pr[M(D)\in A]\leq e^\epsilon\Pr[M(D')\in A]+\delta.
 \label{eq:fed-dp-definition-rewrite}
\end{equation}
相邻关系必须先定义。样本级保护隐藏一条记录的影响，客户端级保护隐藏一个客户端的全部贡献。样本级裁剪不能直接宣称客户端级隐私。

以客户端级加入或删除邻接为例，令上传差分 $u_k$ 裁剪为
\begin{equation}
 \bar u_k=u_k\min\{1,C/\|u_k\|_2\}.
\end{equation}
若固定聚合规模 $m$ 并加入高斯噪声，可以发布
\begin{equation}
 M=\frac1m\left(\sum_k\bar u_k+\xi\right),
 \qquad \xi\sim\mathcal N(0,\sigma^2C^2I_d).
 \label{eq:fed-client-dp-rewrite}
\end{equation}
$C$ 是单个客户端贡献的范数上限，$d$ 是参数维度。裁剪减少敏感度但会产生偏差。噪声标准差、参与客户端数和维度共同影响效用。多轮发布还要进行隐私会计，累计预算不能由单轮 $(\epsilon,\delta)$ 直接替代。

\section{评价协议与失败情形}

联邦实验首先要固定总体目标。若按样本量加权，报告总体样本风险。若关心工厂平等，应报告客户端宏平均和最差客户端指标。分类任务至少报告整体准确率、各客户端召回率和少数类召回率。平均准确率高而某个小工厂完全漏报时，系统仍可能不合格。

每次实验应记录客户端总数、每轮参与数、抽样规则、有效样本量、本地步数、上下行字节数、端侧时间、掉线率、服务器时间和隐私预算。比较方法时固定总通信预算或总端侧计算预算。只固定通信轮数会让本地步数不同的方法使用不同计算量。

测试集应按部署条件隔离。已参与工厂的未来时间段测试旧客户端适应能力。完全未参与的工厂测试跨客户端泛化。新工厂若需要本地适应，应另外划分适应集和最终测试集。测试标签不能参与客户端抽样、超参数选择或停止规则。

以下情形经常使联邦方案失效。客户端标签定义不同却直接平均分类头。网络掉线与故障类型相关却仍把上线客户端当作随机样本。小客户端的少数故障被大客户端的正常样本淹没。模型更新在安全聚合前被服务器读取。差分隐私噪声使小客户端召回下降却只报告总体准确率。恶意客户端利用平均聚合植入后门。模型和归一化统计量版本不一致导致更新坐标失配。基础模型的适配器虽小，却仍包含敏感任务信息。

对多工厂轴承诊断，可以构造三类验证。第一类改变故障比例，检查 label skew。第二类让工厂只观测部分类别，检查类别缺失。第三类改变转速、负载或传感器，检查 feature shift。对每类情形比较 FedAvg、FedProx、SCAFFOLD 和个性化头，并保持相同预算。随后加入掉线、更新反演和恶意客户端测试。只有平均性能、最差客户端、通信成本和威胁模型都满足要求时，才能把方案交给现场验证。

\section*{辅助学习材料}
\begingroup
\emergencystretch=3em
\begin{itemize}
\item McMahan 等人的 FedAvg 论文，说明通信轮次与本地更新的基本折中。\href{https://arxiv.org/abs/1602.05629}{原文}。
\item Kairouz 等人的综述，系统总结统计异质性、系统约束、隐私和攻击面。\href{https://arxiv.org/abs/1912.04977}{原文}。
\item Li 等人的综述，讨论联邦学习的优化、系统和隐私问题。\href{https://arxiv.org/abs/2007.07213}{原文}。
\item Yao 等人的综述，题为 ``Federated Large Language Models: Current Progress and Future Directions''，讨论联邦大语言模型的训练和参数高效适配。\href{https://arxiv.org/abs/2409.15723}{原文}。
\item Jiang 等人的综述，题为 ``Federated Large Language Models: Feasibility, Robustness, Security and Future Directions''，集中讨论可行性、鲁棒性和安全边界。\href{https://arxiv.org/abs/2505.08830}{原文}。
\item Yang 等人的综述，题为 ``Federated Prompt Learning: A Unified Framework, Empirical Analysis, and Future Directions''，讨论提示参数共享、个性化和评测。\href{https://arxiv.org/abs/2608.13844}{原文}。
\item Nguyen 等人的综述，题为 ``From Data Heterogeneity to Convergence: A Data-Centric Review of Federated Learning''，从数据异质性分析训练收敛问题。\href{https://arxiv.org/abs/2606.10595}{原文}。
\item Bonawitz 等人的安全聚合论文，说明掉线和威胁模型下的协议设计。\href{https://arxiv.org/abs/1611.04482}{原文}。
\item TensorFlow Privacy 文档可用于复习裁剪、噪声和隐私会计。\href{https://tensorflow.google.cn/responsible_ai/privacy}{在线文档}。
\end{itemize}
\endgroup

\paragraph{本章小结}
联邦学习把原始数据分散保存时的协作训练写成一个带通信和信任边界的优化问题。客户端权重决定总体目标，参与概率决定抽样偏差，两者不能混用。FedAvg 在一步全参与的特殊条件下等价于集中式梯度下降。多步本地训练会在 non-IID 数据上产生客户端漂移。FedProx 用近端项限制偏离，SCAFFOLD 用控制变量校正方向，FedOpt 在服务器端重新优化聚合差分。个性化方法允许共享规律和本地规律分开保存。

安全聚合限制服务器看到单个更新，差分隐私限制相邻数据集的输出分布差异，鲁棒聚合处理异常贡献。三者的保证对象不同。数据不出域不能单独推出隐私、无投毒或公平性结论。实际验收还要覆盖掉线、少数客户端、未见工厂、通信成本和攻击测试。

联邦学习解决数据协作边界，不能替代任务定义、数据质量和独立评价。下一章讨论流形学习时，仍需检查样本距离是否反映任务相关变化。联邦训练得到的共享表示若把设备身份当作主要差异，低维几何也可能误导后续分析。
